\documentclass[11pt]{article}

\usepackage[T1]{fontenc}
\usepackage[utf8]{inputenc}
\usepackage{amsmath,amssymb}
\usepackage{booktabs}
\usepackage{tabularx}
\usepackage{longtable}
\usepackage{geometry}
\usepackage{hyperref}
\usepackage[numbers]{natbib}
\usepackage{microtype}
\usepackage{tikz}
\usetikzlibrary{arrows.meta,positioning,shapes.geometric,fit}
\geometry{margin=1in}

\newcommand{\RaS}{Repository-as-State}
\newcommand{\RRI}{Repository Resumability Index}
\newcommand{\RTF}{Reconstruction Token Fraction}

\title{Repository-as-State:\\Externalising Agent Continuity for Stateless High-Capability Reasoning}
\author{Paul Maddison\\Independent Researcher\\Liverpool, United Kingdom}
\date{Paper v0.1 --- 2 September 2026}

\begin{document}
\maketitle

\begin{abstract}
Long-horizon software-engineering agents are often framed as persistent intelligent processes whose conversational history, working memory, execution session, or model-side state carries continuity across tasks. This paper investigates an inversion of that architecture. \emph{Repository-as-State} (RaS) treats the software project and its durable repository artefacts as the persistent object while allowing high-capability reasoning processes to be transactional and disposable. At step $t$, a reconstruction function $\Gamma$ selects a bounded reasoning view $X_t$ from durable repository state $R_t$ and current task $U_t$; a high-capability reasoner proposes a bounded engineering transition; a separate execution environment may apply it; deterministic or behavioural evidence validates the candidate state; and only an accepted state is committed as $R_{t+1}$. The reasoner may then cease to exist. A later fresh reasoner should, if the hypothesis holds, recover the decision-relevant engineering state from $R_{t+1}$ without receiving the complete previous reasoning history.

We formalise approximate repository sufficiency, distinguish repository-as-state from repository-as-prompt, and derive a simple cumulative-context model in which naive full-history replay grows quadratically under explicit assumptions while bounded reconstruction grows linearly. The model is not a claim about universal commercial-agent monetary cost: caching, compaction, retrieval, prompt-prefix reuse, and provider architecture may change realised costs. Reconstruction is therefore treated as a first-class falsification route and measured through tokens, repository reads, searches, tool/model calls, elapsed time, and the proposed Reconstruction Token Fraction. We further propose Forced-State-Reset Evaluation and the Repository Resumability Index (RRI), and separate the continuity hypothesis from a secondary tiered-execution architecture in which high-capability reasoning and operational privilege need not reside in the same component.

Controlled empirical evaluation has not yet been executed. This paper contributes the theoretical architecture, formal models, measurement framework, experimental conditions, security analysis, and preregistration-oriented method that will be used to test whether high-quality reasoning can be made disposable without making engineering continuity disposable.
\end{abstract}

\section{Introduction}
\label{sec:introduction}

Software-engineering agents increasingly operate across repositories, shells, test runners, browsers, issue trackers, build systems, and other development tools. SWE-bench established repository-level issue resolution as a measurable software-engineering task \citep{jimenez2024swebench}; SWE-agent showed that the design of the agent--computer interface materially affects a language model's ability to navigate repositories, edit files, and execute tests \citep{yang2024sweagent}; and OpenHands provides a general platform and later SDK for software agents with sandboxed execution, lifecycle control, memory management, and model routing \citep{wang2024openhands,wang2025openhandsSDK}. These systems demonstrate substantial progress toward agents that act like software developers. They do not settle a different systems question: \emph{what must remain persistent when engineering work extends beyond one reasoning session?}

The conventional framing makes the intelligent developer the persistent object. A long-running agent accumulates conversation, observations, summaries, working memory, execution state, and tool history so that later decisions can depend on earlier reasoning. This is useful and sometimes necessary. Modern agent-memory research explicitly studies how experience should be stored, consolidated, retrieved, reflected upon, and reused \citep{du2026memory,luo2026memory}. Context-engineering practice likewise treats the model-visible token budget as a resource that must be curated over long interactions \citep{anthropic2025context}. RaS does not argue that these mechanisms are universally unnecessary.

It asks whether software engineering permits a stronger inversion:

\begin{quote}
\textbf{The persistent object need not be the developer; it can be the software.}
\end{quote}

Software projects already possess unusually rich domain-native durable state. Source code captures implementation. Tests encode executable behavioural constraints. Schemas and interfaces capture contracts. Configuration and build definitions capture environmental assumptions. Architecture records and documentation can preserve rationale and non-obvious constraints. Git history captures accepted transitions. Continuous-integration evidence can attest that particular states compiled, tested, or satisfied external checks. A software project therefore has a natural substrate into which the consequences of prior reasoning can be externalised.

\emph{Repository-as-State} (RaS) investigates whether that substrate can carry enough engineering continuity that the high-capability reasoning process no longer needs to remain alive between bounded tasks. The central proposition is:

\begin{quote}
\textbf{Durable engineering continuity does not necessarily need to live inside the expensive high-capability reasoning process.}
\end{quote}

The distinction between “cheap” and “disposable” is central. RaS does not primarily attempt to make high-quality reasoning cheap. It attempts to make high-quality reasoning \emph{disposable}. A high-capability reasoning process may reconstruct relevant engineering state, reason deeply, produce a bounded transition or work package, delegate routine execution, interpret deterministic evidence, persist the accepted result into the repository, and then cease to exist. A later completely fresh reasoner should need the updated repository and the new task, rather than the full conversational history of its predecessors.

Figure~\ref{fig:ras-architecture} summarises the proposed loop.

\begin{figure}[htbp]
\centering
\input{figures/ras-architecture}
\caption{Conceptual Repository-as-State architecture. The durable object is repository state; reasoning and execution processes may be replaced between validated transitions. No empirical performance is implied.}
\label{fig:ras-architecture}
\end{figure}

The architecture can be written at a high level as
\[
R_t + U_t
\rightarrow
\text{bounded reconstruction}
\rightarrow
\text{high-capability reasoning}
\rightarrow
\text{execution}
\rightarrow
\text{validation}
\rightarrow
R_{t+1}.
\]
After the state transition, the reasoning process may be destroyed. This turns continuity into a property of the engineering programme rather than a property of one long-lived model session.

The idea arose from repeated practical software-engineering workflows in which difficult reasoning was performed in high-capability sessions, detailed implementation packages were handed to comparatively shallow execution agents, and compilers, tests, runtime observations, documentation, and Git persisted accepted results. Fresh reasoning sessions could often resume from repository state. That observation is valuable motivation because the architecture was recognised through use rather than invented purely as an abstract construction. It is nevertheless \emph{observed motivation}, not controlled empirical proof.

RaS therefore frames seven contributions at paper v0.1:

\begin{enumerate}
    \item a Repository-as-State architecture in which durable engineering state is primary and high-capability reasoning is potentially ephemeral;
    \item a theoretical formulation of approximate repository sufficiency;
    \item an explicit separation between durable repository state and the bounded reasoning view reconstructed for a task;
    \item a cumulative-context and provider-neutral cost model under stated assumptions, with reconstruction cost treated as a central falsifier;
    \item Forced-State-Reset Evaluation as a methodology for deliberately removing conversational continuity;
    \item the proposed Repository Resumability Index (RRI) and a separate state-reconstruction probe for analysing resumption;
    \item a secondary tiered-execution and least-privilege architecture that separates reasoning capability from routine operational permissions.
\end{enumerate}

The paper does not claim empirical findings. P0, the initial Forced-State-Reset Pilot, has not been executed. The objective of v0.1 is to define the theory and evaluation clearly enough that subsequent evidence can strengthen, weaken, narrow, or falsify the hypotheses without retroactively changing what was meant by RaS.

\section{The Agent Continuity Problem}
\label{sec:continuity-problem}

Long-running software-engineering work is path dependent. A task performed today may depend on an architectural decision made weeks earlier, a regression discovered in a previous branch, an interface guarantee encoded by another module, a deployment assumption, or a rejected approach whose failure was expensive to understand. A future actor must recover enough of that state to continue correctly.

Agentic engineering systems can accumulate many forms of transient state: conversational history, tool observations, repository exploration, build results, hypotheses, rejected approaches, architectural decisions, environment state, retry history, model outputs, and task progress. Conventional long-running agents may preserve some combination of persistent conversation, compacted context, summaries, episodic memory, working memory, server-side state, execution sandboxes, cached prompts, and persistent tool or session state.

These mechanisms can be valuable. Memory systems such as MemGPT explicitly manage information across memory tiers to extend effective interaction beyond one context window \citep{packer2023memgpt}. Recent surveys describe increasingly sophisticated agent-memory mechanisms spanning storage, reflection, abstraction, retrieval, and policy-controlled management \citep{du2026memory,luo2026memory}. RaS does not claim that such mechanisms are universally unnecessary or inferior.

The narrower question is:

\begin{quote}
\textbf{How much of software-engineering continuity can instead be recovered from authoritative domain-native state?}
\end{quote}

Software engineering is unusual because its domain-native state is rich and executable. Relevant artefacts can include:
\begin{itemize}
    \item source code and generated interfaces;
    \item tests and hidden behavioural contracts;
    \item Git history and accepted diffs;
    \item configuration and toolchain definitions;
    \item schemas, migrations, and compatibility declarations;
    \item architecture decisions and security boundaries;
    \item product or behavioural documentation;
    \item CI evidence and reproducible build information.
\end{itemize}

Let $H_t$ denote accumulated interaction history, $R_t$ durable repository state, $U_t$ the current task, and $S$ a successful engineering outcome. A persistent architecture can condition on:
\begin{equation}
P(S\mid H_t,R_t,U_t).
\end{equation}
The RaS continuity hypothesis asks whether, for a useful class of long-horizon engineering tasks,
\begin{equation}
P(S\mid R_t,U_t)
\approx
P(S\mid H_t,R_t,U_t).
\label{eq:continuity}
\end{equation}

Equation~\ref{eq:continuity} is an empirical target, not an asserted identity. If removal of $H_t$ materially reduces success after controlling for model, repository state, task, and execution resources, then repository state is not sufficiently carrying continuity for the tested workflow.

The hypothesis is plausible because accepted engineering work is routinely reified. Requirements become interfaces and tests. Architectural choices become module boundaries and decision records. Bugs become fixes and regressions tests. Dependency assumptions become lockfiles and compatibility checks. This conversion means that a repository can preserve consequences of earlier reasoning even when it does not preserve the reasoning trace.

Yet important information can remain outside the repository. A conversation may carry stakeholder nuance, an operational caveat, or knowledge that an apparently reasonable alternative was rejected for a subtle reason. Documentation may be stale. Tests may encode only happy paths. Build state may depend on local conditions. External systems may be authoritative for runtime facts. RaS therefore cannot equate “version controlled” with “sufficient”.

There is also an authority problem. Conversation history contains plans, guesses, abandoned hypotheses, and stale observations. Repository artefacts can also be wrong, but many are closer to the engineered object and can be independently checked. A compiler can reject a state. A test can demonstrate a violated invariant. A schema can constrain an interface. A commit can be reviewed. This motivates the systems principle:
\begin{quote}
\textbf{Persist authoritative state. Reconstruct computation.}
\end{quote}

The principle is intentionally limited at this stage to software engineering. Whether analogous authoritative domain state can support disposable reasoning in other domains is future work.

\section{Repository-as-State}
\label{sec:ras}

Repository-as-State models an engineering programme as a sequence of validated repository transitions. At time $t$, the durable repository is $R_t$ and the current task is $U_t$. A reasoning transaction produces a proposed transition, an execution process applies bounded changes, deterministic and hidden checks adjudicate the result, and accepted state becomes $R_{t+1}$:
\begin{equation}
(R_t, U_t)
\rightarrow
\mathcal{Q}_t
\rightarrow
\Delta_t
\rightarrow
\mathcal{E}_t
\rightarrow
\mathcal{V}_t
\rightarrow
R_{t+1},
\end{equation}
where $\mathcal{Q}_t$ is high-capability reasoning, $\Delta_t$ a bounded work package, $\mathcal{E}_t$ execution, and $\mathcal{V}_t$ validation.

The architecture has five important properties.

\subsection{Repository state is durable; reasoning state is not}

The accepted output of a transaction must be represented in artefacts that survive the session: source, tests, build rules, schemas, documentation, history, or evidence. The reasoning process is allowed to disappear completely after the transition. This is stronger than periodically summarising a long-lived agent. The reset condition is architectural, not accidental.

\subsection{Repository-as-State is not repository-as-prompt}

A repository can contain millions of tokens and an arbitrarily long Git history. RaS does not propose placing all of that content into every model call. Durable state may be large; active inference context should remain selective.

Let $K_t$ denote the reconstructed context supplied from $R_t$ for task $U_t$. RaS requires a reconstruction policy
\begin{equation}
\mathcal{R}(R_t,U_t) \rightarrow K_t
\end{equation}
that attempts to preserve task-relevant information while excluding irrelevant repository mass. Reconstruction may involve file search, symbol tracing, test discovery, recent history, architecture records, static analysis, or other repository-local tools. The research question is whether the cost and error rate of this reconstruction remain acceptable as repositories and task depth grow.

\subsection{Continuity is encoded in heterogeneous artefacts}

Source code alone is unlikely to be sufficient. Consider a duplicate-command bug. Historical reasoning may establish the invariant “duplicate requests must not create duplicate effects”. Once captured by an idempotency test, future actors can recover an important semantic constraint without replaying the original debugging conversation. An architecture decision record may preserve why two modules are separated; a schema may preserve a compatibility contract; a lockfile may preserve an environmental assumption.

Condition C in the experimental design therefore performs semantic-state ablation. It retains source, tests, and build rules while removing selected semantic artefacts according to a preregistered manifest. If resumability falls, the experiment begins to identify which repository components actually carry continuity.

\subsection{Accepted state transitions require adjudication}

A language model's statement that a task is complete is not sufficient evidence. RaS places compilers, tests, runtime measurements, hidden behavioural verifiers, and where necessary preregistered human review outside the reasoning process. The repository becomes authoritative only to the extent that its transitions are checked.

This design also reduces dependence on a particular model's self-consistency. The same repository invariant can be read by different models and enforced by model-independent tests. That property is central to the proposed forced-reset experiments.

\subsection{The programme, not the session, is the persistent object}

A long-lived agent architecture naturally treats the agent process as the continuing entity. RaS instead treats the engineering programme as a chain of repository states:
\begin{equation}
R_0 \rightarrow R_1 \rightarrow R_2 \rightarrow \cdots \rightarrow R_n.
\end{equation}
Each transition may be produced by a different fresh reasoning process. Continuity succeeds if those processes can reconstruct enough state to continue correctly.

This perspective does not deny the value of memory systems. It changes which state is considered primary. Agent memory may remain useful for convenience, latency reduction, or personalisation, but the core engineering programme should remain recoverable from durable authoritative artefacts when conversational state is unavailable.

The strongest form of RaS would show high resumability, bounded reconstruction cost, and comparable engineering quality across deep task sequences. A weaker but still useful result could show that RaS works only for particular task localities, repository organisations, or classes of artefact. Failure would also be informative: it would identify information that repositories systematically fail to preserve.

\section{Externalising Engineering Continuity}
\label{sec:externalising}

The engineering state required for correct continuation is broader than source code. A long-running software task accumulates observations, failed hypotheses, accepted constraints, test outcomes, environment assumptions, architectural decisions, migration state, and knowledge of what has already been tried. A persistent agent can carry some of this information in its history. RaS asks which parts should instead be converted into domain-native durable state.

The transformation can be viewed as a form of \emph{decision-relevant semantic compression}. An engineering episode might contain several incorrect hypotheses, repeated debugging attempts, discarded implementations, and long discussion. Once the problem is resolved, future correctness may depend only on a smaller set of durable consequences:
\[
\text{transient reasoning trajectory}
\longrightarrow
\left\{
\begin{array}{l}
\text{accepted implementation},\\
\text{regression test},\\
\text{concise rationale},\\
\text{schema or interface update},\\
\text{explicit operational constraint}.
\end{array}
\right.
\]

The goal is not lossless compression of conversation. Preserving every intermediate thought would defeat the purpose and could preserve noise, stale hypotheses, or model-specific reasoning traces. The goal is to retain information whose absence would change a future correct engineering decision.

This distinction explains why repository state can be semantically richer than a simple code snapshot. Consider a reliability bug in which duplicate requests can create duplicate effects. The original debugging conversation may be lengthy. The durable continuity requirement is narrower: future implementations must preserve idempotency. A regression test, an interface contract, and a concise design note may carry that knowledge more reliably than a conversational summary because they are versioned, reviewable, and partly executable.

Semantic compression is not automatically safe. Final code may reveal what was done without explaining why. Tests may encode only a subset of intended behaviour. Comments can be stale. Architecture records can be wrong. A rejected alternative may matter later because a future engineer could rediscover and repeat the same dangerous approach. RaS therefore predicts that useful continuity will be distributed across heterogeneous artefacts rather than concentrated in one canonical “memory file”.

The architecture also creates pressure toward explicitness. A persistent conversation can carry undocumented assumptions for many turns. A disposable reasoning process cannot. If a constraint matters beyond the current transaction, it should be represented in a durable form appropriate to the domain. That may improve inspectability, but it can also create documentation overhead. The experimental programme must therefore measure whether the required externalisation burden remains practical.

The distinction between authoritative domain state and model-native memory is central. A conversation may contain a statement that an API is idempotent; a test can demonstrate the property against the implementation. A model memory may remember a migration decision; a schema and migration manifest can make the current state inspectable. A session may remember that a compiler version is required; a toolchain file can make the constraint reproducible. In each case, the repository representation is closer to the engineered object and can be consumed by different humans, models, and deterministic systems.

This does not imply that all valuable engineering knowledge belongs in Git. Runtime incidents, production telemetry, external tickets, security systems, or customer requirements may live elsewhere. The narrow RaS claim concerns a software-engineering programme whose authoritative continuity is substantially reconstructable from repository state plus the explicitly supplied current task. Later work may extend the concept to a broader set of authoritative domain stores, but that is a separate hypothesis.

The initial experiment therefore treats externalisation as both a design principle and a measurable variable. If forced-reset agents fail because crucial context was never encoded durably, the result reveals a continuity gap. If semantic-state ablation removes architecture notes or state documents and success falls, that is evidence that descriptive artefacts carry information not recoverable from code and tests alone. If ablation has no effect, the same information may be redundant elsewhere.

RaS should therefore not be evaluated by asking whether an agent can “remember” a repository. The meaningful question is whether an engineering programme can survive the destruction of the reasoning process because enough decision-relevant state has been externalised into durable, reconstructable, independently verifiable artefacts.

\section{From Stateful Agents to Transactional Reasoning}
\label{sec:transactional}

The operational unit in RaS is a reasoning transaction:
\begin{equation}
R_t + U_t
\rightarrow
\text{high-capability reasoning transaction}
\rightarrow
\text{validated state transition}
\rightarrow
R_{t+1}.
\end{equation}
After $R_{t+1}$ is accepted, the reasoning state may be destroyed.

The term \emph{transactional} is used deliberately but cautiously. Database transactions provide formal isolation and atomicity properties that a model-driven engineering workflow does not automatically possess. The analogy concerns lifecycle and durability: a bounded computation reads authoritative state, attempts a change, and only validated output becomes durable programme state.

This reframing has several consequences.

First, transaction boundaries force engineering decisions to become explicit. A persistent conversation can carry unstated assumptions for many turns. A disposable session cannot rely on that hidden continuity. If an assumption will matter later, the architecture pressures the workflow to encode it in source, tests, schemas, documentation, issue metadata, or another durable artefact. This can be viewed as a form of semantic normalisation: ephemeral reasoning is converted into externally inspectable engineering state.

Second, reasoning transactions can be independently scheduled. A programme with many unrelated tasks need not keep one high-capability process resident for each repository. A scheduler could create a fresh reasoner when a task reaches a difficult decision point, provide reconstructed state, collect a bounded output, and release the reasoning resource. This is an infrastructure implication rather than an empirically established saving.

Third, transaction boundaries make reset experiments tractable. Persistent agents and stateless requests can otherwise differ along many hidden dimensions. RaS defines an explicit intervention: after each eligible task, destroy the reasoning session and prohibit resume, previous conversation, external summaries, and copied hidden memory. The next process receives only the allowed repository state and current task. The resulting success or failure is evidence about where continuity resides.

Fourth, transactions permit differentiated failure handling. If an executor cannot compile a proposed change, the repository need not advance. The failure can be returned to the current reasoner, or if the protocol allows, encoded as durable diagnostic evidence for a later transaction. This resembles optimistic state transition more than conversational continuity: invalid proposals are not committed simply because a model produced them.

Fifth, transactional reasoning creates a natural interface for model routing. Different tasks may require different capability tiers. A simple refactor with strong tests may not require the same model as a cross-cutting architectural diagnosis. RaS does not prescribe a routing algorithm, but its bounded transaction interface separates the continuity problem from the choice of model that performs a particular unit of reasoning.

The architecture is therefore compatible with, but distinct from, existing agent frameworks. SWE-agent emphasises the interface through which a model interacts with a computer \citep{yang2024sweagent}. OpenHands exposes composable agent execution, lifecycle, and sandbox mechanisms \citep{wang2024openhands,wang2025openhandsSDK}. RaS asks whether the high-capability reasoning lifecycle itself can become ephemeral while the repository remains the durable programme object.

A major risk is transaction underspecification. If a bounded work package omits a crucial constraint that existed only in conversation, a shallow executor may implement the wrong behaviour. That failure is not an edge case; it directly tests repository sufficiency and decomposition quality. The experimental design should therefore record whether failures arise during reconstruction, reasoning, specification transfer, execution, or verification.

The long-term systems implication is a change in what scales. Instead of scaling a population of persistent long-lived reasoning processes, infrastructure could scale independent reasoning transactions over durable engineering state. Whether that improves cost, throughput, reliability, or security remains a hypothesis to be tested.

\section{Behavioural Repository Sufficiency}
\label{sec:sufficiency}

The original v0.1 draft expressed repository sufficiency through equality of latent distributions over “useful next actions” and a KL-divergence bound. Hostile review rejects that formulation for the planned experiments. The action space is open-ended, equivalent implementations need not have comparable surface actions, and the proposed runs do not expose a well-defined probability distribution over engineering actions. KL divergence would therefore add mathematical appearance without an estimable quantity.

The paper instead defines an observable behavioural target.

Let $Y_{s,d,c}\in\{0,1\}$ denote preregistered hidden-verifier success for task sequence $s$, depth $d$, and condition $c$. Let $c=A$ be the persistent-history control and $c=B$ the forced-reset condition. The two conditions must use the same accepted starting repository state, current task, high-capability model/configuration, tool contract, workspace image, and acceptance criteria. The intended treatment difference is availability of predecessor reasoning-session history.

For later repeated confirmatory experiments, define:
\begin{equation}
p_c(d)=P(Y=1\mid c,d,\mathcal{M},\mathcal{E},\mathcal{T}),
\label{eq:behavioural-success}
\end{equation}
where $\mathcal{M}$ identifies the matched model/configuration, $\mathcal{E}$ the controlled execution environment, and $\mathcal{T}$ the task-sampling population.

A future confirmatory study can preregister a scientifically justified non-inferiority margin $\delta>0$ and ask whether:
\begin{equation}
p_B(d)-p_A(d) \ge -\delta
\label{eq:noninferiority-target}
\end{equation}
over the declared task population and depth range. If symmetric equivalence is scientifically required, a two-sided margin can be preregistered instead. The margin must be fixed before outcomes and justified in engineering terms; this paper does not invent one.

P0 is too small to establish non-inferiority statistically. Its role is to test whether the intervention, task corpus, hidden verification, leakage controls, and telemetry are workable and whether an obvious continuity collapse appears.

“Repository sufficiency” is therefore shorthand for a conditional behavioural claim, not an intrinsic property of a repository. A repository may be adequate for one task class, model, tool interface, and depth but not another. It also need not preserve every past thought. The question is whether accepted durable state permits matched fresh reasoning to achieve comparable engineering outcomes at acceptable reconstruction cost.

Operational evidence includes:
\begin{itemize}
    \item hidden-verifier task success;
    \item complete-sequence success across dependent tasks;
    \item success by sequential depth;
    \item reconstruction-probe accuracy;
    \item regression rate;
    \item human-independent resumption;
    \item reconstruction resource use and elapsed time.
\end{itemize}

Semantic-state ablation remains useful but must test marginal continuity value rather than “making the repository worse”. Condition C should remove one small preregistered artefact class or information channel at a time while keeping task semantics and buildability intact. Redundant encoding must be audited: deleting an ADR proves little if the same decision is reproduced verbatim in a test, comment, or later allowed commit.

The strong RaS hypothesis is thus intentionally narrow: for a declared class of dependent software-engineering tasks at accepted transition boundaries, predecessor reasoning-session history may have sufficiently small marginal behavioural value that a fresh reasoner can continue from authoritative project state without material outcome loss. Whether that is true is an empirical question.

\subsection{What the accepted Subject-B evidence shows}

Subject-B P1 observed equal non-zero behavioural performance: 18/30 behaviours in each condition, with 30 matched agreements and no disagreements across three tasks. Subject-B P2 expanded the same-repository design to four tasks and three repetitions. It observed 0/39 satisfied behaviours in both conditions, 39 matched agreements, no disagreements, and 0/12 candidate-level overall passes in each condition. The P2 result is therefore a floor effect: it detects no A/B difference under this design but does not demonstrate successful behavioural preservation or repository-state sufficiency. Neither result supports an equivalence or non-inferiority claim, and the two studies must not be pooled into one correctness score.

\section{Context Growth, Reconstruction Cost, and State Economics}
\label{sec:cost}

RaS has an economic motivation, but its central intellectual claim is continuity placement rather than a specific price forecast. A provider-neutral model is therefore preferable to assumptions about internal GPU scheduling, proprietary caches, margins, or subscription economics.

\subsection{A simple cumulative-context model}

Let $B$ be fixed base context, $g$ the average number of new historical tokens introduced per sequential interaction, and $n$ the number of reasoning steps. Under a naive full-history replay model, the input presented at step $i$ is
\begin{equation}
T_i = B + g(i-1).
\end{equation}
The cumulative supplied context is
\begin{align}
T_{\mathrm{persistent}}(n)
&=
\sum_{i=1}^{n}[B+g(i-1)] \\
&=
nB + \frac{gn(n-1)}{2}.
\label{eq:persistent-context}
\end{align}
For fixed $g>0$,
\begin{equation}
T_{\mathrm{persistent}}(n)\in\Theta(n^2).
\end{equation}

Now let $K_i$ be the repository/task context reconstructed for RaS step $i$. If
\begin{equation}
K_i \le K
\end{equation}
for bounded $K$, then
\begin{equation}
T_{\mathrm{RaS}}(n)
=
\sum_{i=1}^{n}K_i
\le nK,
\end{equation}
and, given a non-degenerate lower bound, cumulative supplied context is $\Theta(n)$.

Figure~\ref{fig:context-growth} illustrates only these mathematical assumptions.

\begin{figure}[htbp]
\centering
\input{figures/context-growth}
\caption{Conceptual cumulative-context growth under the stated models. The curves are illustrative mathematical shapes, not measured agent or provider costs.}
\label{fig:context-growth}
\end{figure}

This result must not be confused with a claim that commercial persistent agents universally have quadratic monetary cost. Real systems can use prompt-prefix caching, compaction, summarisation, context pruning, retrieval, server-side state, and changing provider architectures. Agent interfaces can also avoid replaying observations that are no longer useful. The theorem concerns cumulative information supplied under a specified naive replay model.

The bounded-$K$ assumption is equally important and potentially false. A more realistic model is:
\begin{equation}
K_i =
f(
|R_i|,
L_i,
D_i,
Q^{\mathrm{doc}}_i,
Q^{\mathrm{test}}_i,
Q^{\mathrm{arch}}_i,
Q^{\mathrm{retrieval}}_i
),
\label{eq:realistic-k}
\end{equation}
where $|R_i|$ is repository scale, $L_i$ task locality, $D_i$ dependency width, and the $Q$ terms represent documentation, test, architecture, and retrieval quality. The empirical question is not whether $K_i$ is literally constant. It is whether $K_i$ grows slowly enough, and with sufficiently high reconstruction fidelity, for externalised continuity to remain useful.

\subsection{Reconstruction cost as a falsification route}

The strongest obvious objection is that RaS merely moves cost from persistent context into repeated repository rediscovery. This is legitimate. RaS therefore defines reconstruction cost $C_{\mathrm{reconstruction}}$ as a first-class outcome rather than an implementation detail.

Future experiments should measure, where available:
\begin{itemize}
    \item input tokens attributable to reconstruction;
    \item files and bytes read;
    \item repository and symbol searches;
    \item history and dependency-navigation operations;
    \item tool calls and model calls;
    \item elapsed reconstruction time;
    \item cached versus uncached context;
    \item retries caused by incomplete or incorrect reconstruction.
\end{itemize}

We propose the \emph{Reconstruction Token Fraction}:
\begin{equation}
\mathrm{RTF}
=
\frac{T_{\mathrm{reconstruction}}}
{T_{\mathrm{RaS,input}}}.
\label{eq:rtf}
\end{equation}
RTF is meaningful only with a consistent classification rule. The boundary between “reconstruction” and “task reasoning” must be preregistered before outcomes are observed. Otherwise repository reading can be relabelled after the fact to make the architecture appear cheaper.

If reconstruction becomes excessively expensive as repositories grow, that weakens the RaS state-economics thesis even if forced-reset engineering success remains high.

\subsection{Provider-neutral cost decomposition}

For a persistent coding agent:
\begin{align}
C_{\mathrm{persistent}}
=&\ C_{\mathrm{reasoning}}
+ C_{\mathrm{context/history}}
+ C_{\mathrm{agent-state}} \nonumber\\
&+ C_{\mathrm{execution}}
+ C_{\mathrm{tooling}}
+ C_{\mathrm{orchestration}}.
\label{eq:cpersistent}
\end{align}

For RaS:
\begin{align}
C_{\mathrm{RaS}}
=&\ C_{\mathrm{reconstruction}}
+ C_{\mathrm{reasoning}}
+ C_{\mathrm{shallow-execution}} \nonumber\\
&+ C_{\mathrm{repository-state}}
+ C_{\mathrm{orchestration,RaS}}.
\label{eq:cras}
\end{align}

No component is assumed to be zero. The economic hypothesis is that for sufficiently long engineering programmes,
\begin{equation}
C_{\mathrm{RaS}} < C_{\mathrm{persistent}}
\end{equation}
while
\begin{equation}
Q_{\mathrm{engineering,RaS}}
\approx
Q_{\mathrm{engineering,persistent}}.
\end{equation}
The second condition is essential. Lower nominal model usage is not a saving if regressions, retries, or failed tasks increase sufficiently.

Three categories must remain separate:

\paragraph{Measured cost proxy.}
Observable model tokens, provider-reported usage, calls, files read, searches, elapsed time, execution operations, retries, and session persistence.

\paragraph{Theoretical infrastructure implication.}
A consequence derived from an explicit resource model, such as the possibility that lower persistent state per useful programme permits greater concurrency.

\paragraph{Unobservable provider-internal cost.}
GPU allocation, exact KV-cache cost, storage amortisation, scheduler behaviour, internal cache hit rates not exposed to the user, and provider margin.

The experiment cannot infer the third category from subscription pricing. Nor should a bundled chat product be described as “free inference”.

Table~\ref{tab:persistent-vs-ras} summarises the architectural comparison.

\input{tables/persistent-vs-ras}

Long-context serving research provides evidence that inference-state management can materially affect throughput. PagedAttention, for example, addresses KV-cache fragmentation and sharing in LLM serving \citep{kwon2023vllm}. That result motivates care around state burden but does not establish RaS savings. The defensible RaS claim remains conditional: if less high-capability historical state must persist per useful engineering programme, fixed inference infrastructure may potentially support more useful concurrent work.

\section{Tiered Execution}
\label{sec:tiered}

RaS continuity and tiered execution are related but separable claims. The continuity hypothesis can be tested using the same high-capability model for reasoning and editing. Tiered execution introduces a second question: after high-capability reasoning has produced a bounded work package, can routine stateful execution be delegated to a lower-cost, local, or more constrained worker without materially reducing engineering quality?

The high-capability reasoner is best reserved, under the proposed architecture, for operations such as architecture, decomposition, difficult causal reasoning, complex debugging, review, difficult trade-off analysis, interpretation of evidence, and work-package creation. An execution worker may instead perform repository edits, compilation, test execution, formatting, linting, Git operations, local-environment interaction, and tightly constrained environment actions where permitted.

\input{tables/reasoner-executor}

Let $N_R$ be the number of high-capability reasoning operations, $N_E$ the number of execution operations, $c_R$ the average cost of a high-capability operation, $c_E$ the average cost of an execution operation, and $C_O$ orchestration overhead. A simplified monolithic model is:
\begin{equation}
C_M=(N_R+N_E)c_R.
\end{equation}
Tiered execution is:
\begin{equation}
C_T=N_Rc_R+N_Ec_E+C_O.
\end{equation}
It is cheaper under this theoretical model when:
\begin{equation}
N_E(c_R-c_E)>C_O.
\label{eq:tiered}
\end{equation}

Equation~\ref{eq:tiered} is theoretical. It does not assume that $c_E<c_R$ in all deployments, that execution quality is equivalent, or that orchestration is cheap. A nominally inexpensive executor can become more expensive in outcome-normalised terms if it repeatedly misunderstands the work package, introduces regressions, exhausts retries, or forces the high-capability reasoner to re-enter.

Tiered execution can therefore fail through at least four mechanisms. First, the work package may omit a crucial semantic constraint. Second, the executor may lack enough capability to handle an unexpected repository state. Third, the executor may have insufficient tool access to reproduce the required evidence. Fourth, repeated hand-offs may create latency and context overhead larger than any saving.

Condition D in the broader research programme evaluates this architecture. P0 may initially focus primarily on persistent control, forced reset, and semantic-state ablation because those conditions test the central continuity proposition more directly. Tiered execution should not be allowed to confound the first continuity result.

The architecture nevertheless matters because it highlights a systems principle: the component that performs the most sophisticated reasoning does not necessarily need to perform the most machine operations. If continuity is external, reasoning capability, execution state, and operational privilege can be independently scoped.

\section{Tests as Executable Semantic Memory}
\label{sec:executable-memory}

Software repositories contain artefacts that do more than describe state. Tests can actively enforce knowledge produced by earlier reasoning. This makes them a candidate mechanism for preserving decision-relevant semantic history across agent resets.

Suppose a debugging episode establishes the requirement:
\begin{quote}
Duplicate requests must not create duplicate effects.
\end{quote}
A persistent agent could remember that statement in conversation or explicit memory. A repository can instead retain a regression test such as \emph{duplicate request is idempotent}. Future implementations that violate the invariant are then rejected even if the actor never sees the original conversation.

Tests possess several properties that are attractive for continuity:
\begin{itemize}
    \item they are executable;
    \item they are versioned with the code they constrain;
    \item they can be evaluated independently of the current reasoning model;
    \item they are automatically checkable in CI;
    \item they detect regressions rather than merely documenting intent;
    \item they are human reviewable.
\end{itemize}

Calling tests “memory” is metaphorical unless operationalised. A test does not store every reason an invariant exists, and an implementation can overfit a narrow test while violating broader intent. Nevertheless, tests can carry information that changes the probability of future correct action. In the terminology of Section~\ref{sec:sufficiency}, they can reduce the difference between action conditioned on historical interaction and action conditioned only on repository state.

This motivates the distinction between \emph{lossless conversation memory} and \emph{decision-relevant semantic state}. Engineering routinely performs semantic compression:
\[
\begin{array}{c}
\text{many observations, hypotheses, failed fixes, and discussions}\\
\downarrow\\
\text{accepted fix} + \text{regression test} + \text{optional design record}.
\end{array}
\]
The final repository may be much smaller than the reasoning trajectory yet preserve what matters for future correctness.

This compression can be lossy in harmful ways. Tests may omit edge cases; comments can become stale; design rationale may disappear; source code can encode \emph{what} without \emph{why}. RaS therefore does not claim that tests alone are sufficient. Instead, source, tests, architecture records, schemas, build rules, and history form a heterogeneous state representation.

The semantic-state ablation condition can begin to identify relative contribution. For example, removing architecture documents while retaining tests asks whether executable constraints already encode enough of the relevant decision. Conversely, retaining architecture documents while withholding selected non-essential tests in a later experiment could measure how much continuity depends on executable versus descriptive state.

A further implication is model independence. Conversational memory may be coupled to a provider-specific session representation. A regression test can be consumed by any future actor with access to the repository and test runner. This property makes tests particularly compatible with disposable reasoning and cross-model replication.

The empirical challenge is to avoid circular evaluation. If the same visible tests both communicate the task and define success, an agent can pass without reconstructing the broader intended behaviour. Hidden behavioural verification is therefore required for central experiments. Visible tests can carry semantic state; hidden verifiers should adjudicate whether that state supports genuinely correct continuation rather than test-specific imitation.

\section{Durable Documentation and Selective Retention}
\label{sec:durable-docs}

Code and tests are powerful carriers of engineering state, but they do not necessarily preserve the rationale required for a later correct decision. Some information deserves explicit durable semantic representation: why an architectural boundary exists, which alternatives were rejected for reasons that remain valid, what security assumption a component relies upon, which migration phase is currently active, or which business invariant cannot be inferred from implementation alone.

RaS therefore does not advocate “code is documentation”, nor does it advocate dumping full reasoning transcripts into Markdown. Both extremes are problematic. A code-only repository can erase rationale. A transcript-heavy repository can replace one unbounded history problem with another, preserve obsolete hypotheses, and make reconstruction expensive.

A useful retention principle can be expressed conceptually. Persist an item of engineering knowledge $x$ when:
\begin{equation}
P_{\mathrm{reuse}}(x)
\times
L_{\mathrm{missing}}(x)
>
C_{\mathrm{persist+retrieve}}(x).
\label{eq:retention}
\end{equation}
Here $P_{\mathrm{reuse}}(x)$ is the probability that future work will need the knowledge, $L_{\mathrm{missing}}(x)$ is the expected loss if it is absent, and $C_{\mathrm{persist+retrieve}}(x)$ is the cost of preserving, maintaining, and later reconstructing it.

Equation~\ref{eq:retention} is not proposed as a production formula with directly measurable probabilities and monetary values. It is a design heuristic that makes the trade-off explicit. High-consequence, non-obvious, likely-to-recur knowledge should be externalised. Low-value transient reasoning should usually disappear.

This suggests several durable documentation classes:

\begin{itemize}
    \item \textbf{architecture decisions}: why a boundary, dependency, protocol, or data model exists;
    \item \textbf{operational constraints}: assumptions about ordering, retries, idempotency, time, storage, or failure recovery;
    \item \textbf{security boundaries}: trust assumptions, privilege constraints, data handling, and prohibited flows;
    \item \textbf{migration state}: what has moved, what remains dual-run, and which compatibility guarantees still apply;
    \item \textbf{business intent}: non-obvious requirements that tests may only partially encode;
    \item \textbf{known rejected alternatives}: only where repeating the alternative would create meaningful cost or risk.
\end{itemize}

The quality of these artefacts is a threat to validity as well as a design variable. An unusually disciplined repository may make RaS appear stronger than typical software projects. Conversely, a repository with stale or contradictory documentation may make forced reset fail even when persistent history would have repaired the ambiguity. Semantic-state ablation is therefore not merely a convenience; it helps identify whether success depends on explicit documentation, tests, source structure, history, or redundancy between them.

Repository-local agent instruction files are one contemporary example of explicit semantic state. A 2026 empirical study of more than twelve thousand Cursor rule files found that such files commonly encode project context, engineering practices, structure, configuration, and maintainability guidance \citep{sun2026cursorrules}. RaS does not claim those files are sufficient or ideal. Their existence does, however, show that developers already externalise agent-relevant context into versioned repository artefacts.

The desired end state is a repository that preserves enough high-value semantic information for future correct action without becoming a transcript archive. That balance is itself part of the reconstruction-cost question.

\section{Security and Least Privilege}
\label{sec:security}

Separating durable continuity from the reasoning process has a security implication independent of economic savings: high intelligence does not imply high privilege.

A high-capability reasoner may require repository read access, reasoning capability, and permission to propose a patch or bounded work package. An executor may require workspace write access, compiler and test runner access, a formatter or linter, restricted Git, and tightly scoped environment interaction. Production credentials may belong to neither.

This separation can reduce blast radius only if it is enforced outside the model. Prompt instructions are not an access-control boundary. Repository permissions, sandbox configuration, network policy, credential brokers, command allowlists, approval gates, and deterministic validators must define what each component can actually do.

A compromised or simply mistaken reasoner should not automatically gain production credentials merely because it can perform difficult technical analysis. A compromised executor should not automatically inherit the model's broader contextual access. The architecture can therefore narrow the trusted computing base by making reasoner, executor, validator, and credential authority separately replaceable components.

The separation introduces its own risks.

\paragraph{Compromised reasoning process.}
A reasoner may generate a technically plausible but dangerous work package, exploit a downstream executor, or misclassify untrusted repository content as control instructions. The executor must validate requested operations against policy rather than blindly treating the proposal as authority.

\paragraph{Compromised executor.}
An executor can alter repository state, falsify local evidence, or attempt privilege escalation. Validation should therefore include independently collected evidence where possible, and sensitive transitions should require stronger trust boundaries than ordinary compilation.

\paragraph{Repository poisoning.}
Repositories can contain malicious comments, instruction files, generated text, test fixtures, or documents intended to manipulate an agent. Reconstruction must treat repository content as data with provenance, not as an automatically trusted instruction hierarchy.

\paragraph{Untrusted generated state.}
Agent-generated documentation can itself become durable attack surface. A malicious instruction committed today may be consumed by a fresh model tomorrow. Review and policy should therefore apply to semantic state as well as executable code.

\paragraph{Human approval boundaries.}
High-consequence operations may require explicit human approval even when the reasoner and executor agree. The architecture should make such boundaries visible rather than hiding them inside an autonomous session.

Potential benefits include reduced blast radius, narrower credential distribution, model-independent enforcement, and shorter-lived capability tokens. These are architectural implications, not measured security results. Future work should evaluate policy violations and adversarial repository inputs directly.

The key point is that persistent engineering continuity does not require a persistent highly privileged intelligence. If repository state carries continuity, reasoning lifetime and privilege lifetime can be designed independently.

\section{Repository Resumability and State Reconstruction}
\label{sec:resumability}

A forced reset produces a binary engineering outcome only after the agent has first attempted to understand the repository. RaS therefore separates two questions:

\begin{enumerate}
    \item Can a fresh reasoner reconstruct the relevant engineering state?
    \item Given that reconstructed state, can it perform the next engineering task correctly?
\end{enumerate}

This separation matters because a failed continuation may originate in reconstruction rather than reasoning or execution.

\subsection{Repository Resumability Index}

We propose the \emph{Repository Resumability Index} (RRI):
\begin{equation}
\mathrm{RRI}
=
\frac{
\text{successful correct continuations after complete reasoning-state reset}
}{
\text{eligible forced resets}
}.
\label{eq:rri}
\end{equation}
Conceptually,
\begin{equation}
\mathrm{RRI}
=
P(\text{correct continuation}
\mid
\text{complete conversational reset})
\end{equation}
under a specified experimental protocol.

RRI is a \textbf{proposed research metric}, not an established standard. Its interpretation depends on eligibility and correctness definitions. Compilation alone is insufficient. A successful continuation should satisfy preregistered behavioural acceptance criteria, preferably including hidden verification unavailable as solution guidance.

RRI should be reported by task depth as $\mathrm{RRI}_1,ldots,\mathrm{RRI}_n$, by repository, and where experiments are repeated, with uncertainty intervals appropriate to the design. Infrastructure failures should not automatically count as agent failures, but exclusion must be governed by a predeclared rule. An “eligible forced reset” is one in which the correct starting state, task specification, model configuration, and required experimental infrastructure were available and the reset intervention was successfully enforced.

\subsection{State-reconstruction probe}

Before editing after a reset, the experimental reasoner should produce a structured \emph{state-reconstruction probe}. The probe must not request or record hidden chain-of-thought. It captures concise, inspectable task state:

\begin{itemize}
    \item relevant architecture;
    \item current externally observable behaviour;
    \item affected components;
    \item constraints and invariants;
    \item relevant tests and verification paths;
    \item work apparently completed in prior repository state;
    \item work still required by the current task;
    \item evidence paths supporting the reconstruction;
    \item uncertainty and known unknowns.
\end{itemize}

The probe must be stored \emph{outside} the experimental subject repository. Otherwise it becomes new durable state and contaminates subsequent resets.

Let $\widehat{R}_t$ denote the agent's estimated task-relevant state reconstructed from $R_t$. The probe allows analysis of
\begin{equation}
R_t \rightarrow \widehat{R}_t
\end{equation}
separately from
\begin{equation}
\widehat{R}_t \rightarrow \text{successful engineering transition}.
\end{equation}
A run may reconstruct the repository accurately and still fail implementation; alternatively it may produce a weak reconstruction but succeed through local exploration during execution. These are different mechanisms.

Reconstruction accuracy cannot be measured by comparing private reasoning traces. Instead, a preregistered rubric can compare the probe against a held-out, publication-safe state specification or blinded evaluator labels. The rubric should emphasise constraints, dependencies, prior accepted work, and missing-state errors rather than stylistic similarity.

RRI and reconstruction accuracy are complementary. RRI asks whether the programme continues. The probe asks how well a fresh reasoner understood the state before acting. Reconstruction cost then asks how expensive that understanding was to obtain.

\section{Experimental Method}
\label{sec:method}

The defining experiment for RaS is deliberate destruction of agent continuity. The intervention must be stronger than starting a new visible chat while silently retaining summaries, provider memory, or orchestration state.

For a sequence of engineering tasks, let $A_i$ denote the reasoning process used at stage $i$:
\begin{align}
A_1 + R_0 + U_1 &\rightarrow R_1,\\
\text{destroy }A_1, \qquad
A_2^{\mathrm{fresh}} + R_1 + U_2 &\rightarrow R_2,\\
\text{destroy }A_2, \qquad &\ldots
\end{align}
The reset condition prohibits previous conversation, resume mechanisms, external session summaries, copied hidden memory, and manual transfer of task-specific historical state. Only the allowed repository state and the current task survive.

\subsection{Experimental conditions}

The research programme defines four conditions.

\paragraph{A --- Persistent monolithic control.}
A continuing high-capability agent session retains history and may reason, edit, execute tests, and inspect the environment.

\paragraph{B --- RaS forced reset.}
The same high-capability model is used, but the reasoning session is destroyed after every bounded task. Complete allowed repository state persists. Each new task begins in a fresh reasoning session.

\paragraph{C --- RaS semantic-state ablation.}
Condition C follows B but removes selected semantic repository artefacts according to a preregistered manifest. Candidate removals include architecture decision records, state or session documents, and semantic product documentation. Source, tests, and build rules remain unless a later experiment explicitly studies those artefacts.

\paragraph{D --- RaS tiered execution.}
High-capability reasoning remains transactional while a separate lower-cost, stateful, or local execution worker performs constrained edits, compilation, tests, and Git operations. High-capability reasoning re-enters only under protocol-defined escalation conditions.

\subsection{Task construction}

Sequential tasks must require genuine dependence on earlier accepted state. Independent issue-resolution tasks would not test continuity deeply enough. A task sequence should therefore include evolving interfaces, cross-module invariants, regression-sensitive behaviour, and decisions whose consequences matter several stages later.

Task construction also creates leakage risk. If tasks are derived from historical commits, later repository state can reveal the answer. A valid corpus must pin the starting revision for each stage and exclude future files, tests, comments, and history unavailable at that point. Historical commit survivor bias must be recorded: successful historical changes may underrepresent dead ends and ambiguous requirements.

\subsection{Success and hidden verification}

A task is successful only if it satisfies preregistered acceptance criteria. Visible tests alone are insufficient because they both communicate repository state and can invite test-specific overfitting. Hidden behavioural verification should assess the externally intended behaviour without revealing solution details.

Where hidden verification is impossible, the protocol should use another independent criterion such as a held-out test suite or blinded human assessment. Human intervention must be explicitly limited and recorded; otherwise a human operator can become an unmeasured continuity channel.

\subsection{Reconstruction telemetry}

Condition B is not meaningful without measuring rediscovery. Each transaction should record, where available:
\begin{itemize}
    \item total input and output tokens;
    \item reconstruction-attributed tokens;
    \item cached and uncached input;
    \item files and bytes read before implementation;
    \item search, symbol, and history operations;
    \item model and tool calls;
    \item elapsed reconstruction and total task time;
    \item retries, escalations, and human interventions.
\end{itemize}

Classification rules for what counts as reconstruction must be preregistered. Otherwise an experimenter can make RaS look cheaper by classifying repository-reading work as ordinary reasoning.

\subsection{P0: Forced-State-Reset Pilot}

P0 is a pilot intended to validate the experimental machinery rather than establish generality. The likely private subject is KynticAI Fortress. Proprietary source, private diffs, hidden verifier implementations, and raw traces containing source must remain outside the public repository. Public evidence should consist only of sanitised aggregate results, protocol metadata, and publication-safe artefact hashes.

P0 should test whether reset enforcement is technically credible, whether task sequencing genuinely depends on prior state, whether hidden verification can be kept independent, and whether reconstruction telemetry is complete enough for later analysis. It should also identify variance and failure modes needed for confirmatory power calculations.

\subsection{Analysis}

Primary reporting should include task success by condition and RRI by sequential depth. Secondary outcomes should include regression rate, reconstruction accuracy, RTF, total input tokens, tool calls, model calls, elapsed time, retry rate, and intervention count. Later confirmatory studies should preregister estimands, confidence intervals, missing-data rules, multiplicity control, and power calculations.

The experimental method is designed to permit a negative result. If reset performance degrades substantially, if semantic ablation has no interpretable effect, or if reconstruction cost rises enough to erase resource advantages, those findings directly constrain the RaS claim.

\section{Results}
\label{sec:results}

P0 has not been executed. The controlled results below are the bounded Subject-B P1 and P2 records; neither is a general validation of Repository-as-State.

\subsection{Subject-B P1}

P1 compared persistent-session and fresh-session conditions across three matched tasks. Condition A and Condition B each satisfied 18/30 governing behaviours. The matched vectors had 30 agreements and 0 disagreements; both conditions passed two of the three tasks. This is bounded positive evidence for the narrow P1 observation that fresh reconstruction matched persistent-session correctness in this sample. It is not an equivalence, non-inferiority, universal sufficiency, or resource-cost claim.

\subsection{Subject-B P2 Level 2}

P2 used four tasks, three repetitions, two conditions, 24 executions and 12 matched task-by-repetition pairs. The four task contracts define 13 governing behaviours, giving 39 behaviour observations per condition. Both conditions satisfied 0/39 and failed 39/39. The matched vectors had 39 agreements and 0 disagreements, and candidate-level overall verifier passes were 0/12 in both conditions.

The P2 result is a floor effect. No behavioural correctness difference was observed between the persistent-session and fresh-session conditions in P2, but because both conditions were at the floor, P2 does not demonstrate successful behavioural preservation by either condition and cannot by itself be interpreted as evidence of behavioural equivalence or repository-state sufficiency. P1 and P2 are not pooled.

The first Item-63 adjudication attempt is preserved as invalid history: 24 placeholder calls produced zero valid adjudications, premature unblinding occurred, and one later sealed-verifier diagnostic result was excluded. The accepted recovery used 24 fresh blinded sealed-verifier adjudications, froze them before unblinding, and repaired only a derived denominator error. No P2 task-model or verifier reruns were added.

Resource/reconstruction telemetry remains separate and descriptive; it does not establish a cost or speed advantage.

Future results should report, at minimum, subject or public repository identifier, pinned starting revision, task sequence, condition, model/runtime version, reset depth, hidden-verifier outcome, state-reconstruction score, reconstruction telemetry, total task telemetry, retries, intervention flags, and sanitisation status.

Expected result tables include:
\begin{itemize}
    \item task success and RRI by condition and sequential depth;
    \item state-reconstruction accuracy by condition;
    \item reconstruction tokens, RTF, files read, searches, calls, and elapsed time;
    \item regression and retry rates;
    \item outcome-normalised measured cost proxies.
\end{itemize}

No dummy percentages or illustrative result values are included because plausible numbers could later be mistaken for measurements.

The project uses the claim categories in Table~\ref{tab:claim-categories} to prevent theoretical, observed, and hypothetical statements from being silently presented as findings.

\input{tables/claim-categories}

\section{Related Work and Novelty Boundary}
\label{sec:related}

RaS intersects software-engineering agents, agent memory, context engineering, durable execution, model routing, repository instructions, inference-state management, and repository infrastructure. The contribution is not the invention of any one of those areas; it is the combined continuity architecture and the proposed way of testing it.

\subsection{Software-engineering agents and benchmarks}

SWE-bench introduced a benchmark of real GitHub issues requiring repository-level software changes and execution-based evaluation \citep{jimenez2024swebench}. SWE-agent demonstrated that a purpose-built agent--computer interface can improve navigation, editing, context management, and test execution \citep{yang2024sweagent}. OpenHands broadened this direction into an open platform for generalist software agents with sandboxed execution and benchmark integration \citep{wang2024openhands}. Its Software Agent SDK later made lifecycle control, memory management, model-agnostic routing, and local-to-remote execution explicit architectural concerns \citep{wang2025openhandsSDK}.

These systems establish strong foundations for repository interaction and execution. RaS isolates a different variable: whether high-capability conversational continuity must survive across sequential engineering stages after accepted repository state has been persisted.

\subsection{Agent memory and context engineering}

MemGPT treats extended interaction as a virtual-context management problem, moving information between memory tiers to support long-running conversations and document analysis \citep{packer2023memgpt}. Recent memory surveys organise a broader literature spanning storage, reflection, experience abstraction, retrieval, hierarchical context, and policy-controlled memory \citep{du2026memory,luo2026memory}.

Context engineering addresses the related problem of curating what information enters a finite model context. Anthropic's engineering guidance explicitly recommends maintaining a high-signal context through retrieval, compaction, and selective loading in long-horizon agents \citep{anthropic2025context}. SWE-agent similarly manages environment feedback and command history to avoid unnecessarily verbose context \citep{yang2024sweagent}.

RaS is complementary. Instead of asking only how to preserve and retrieve prior agent experience, it asks how much of the experience can be allowed to disappear after its decision-relevant consequences have been encoded in authoritative software artefacts.

\subsection{Repository-local semantic state}

Repository-local instruction files are a direct example of durable agent-relevant context. Sun et al. analyse more than twelve thousand Cursor rule files and find recurring project context, engineering practices, structure, configuration, maintainability guidance, and some security content \citep{sun2026cursorrules}. RaS does not claim that instruction files are sufficient or uniquely important; they demonstrate that externalising semantic guidance into versioned repository state is already common enough to study empirically.

Tests, schemas, build rules, architecture records, and history extend this idea beyond prompt-like instructions. RaS's semantic-state ablation is intended to identify which classes actually contribute to resumability.

\subsection{Durable execution and compute--storage separation}

Durable workflow systems provide a broader systems precedent for reconstructable volatile compute. Durable Functions formalises a programming model in which execution progress can survive failures under compute--storage separation and record/replay \citep{burckhardt2021durable}. Temporal similarly describes durable workflow execution using persisted event history and replaceable worker processes, although RaS does not depend on Temporal semantics.

RaS differs because language-model reasoning is nondeterministic and its task is to reconstruct semantic engineering state, not replay a deterministic workflow history. The precedent is architectural: volatile computation can operate over durable external state without the worker process itself being the persistent object.

\subsection{Model routing}

RouteLLM studies learned routing between stronger and weaker models to optimise quality/cost trade-offs \citep{ong2024routellm}. Tiered RaS execution may use model routing, but routing is not the continuity contribution. RaS first asks whether the programme can survive destruction of the previous reasoning state; only then does it ask which capability tier should execute a particular transaction.

\subsection{Inference-state management}

PagedAttention and vLLM show that KV-cache allocation and sharing materially affect LLM-serving throughput, particularly for long sequences \citep{kwon2023vllm}. This supports the general systems observation that inference state is a real resource. It does not prove that externalising engineering continuity reduces provider cost or data-centre pressure. RaS keeps those consequences as hypotheses and implications.

\subsection{Repository-state infrastructure}

Cursor's 2026 “Git at any scale” article describes Continuity, a repository system based on a write-ahead log in S3-compatible object storage, with repository-serving state materialised on nodes from durable object state \citep{marti2026continuity}. This is a useful contemporary precedent for separating repository-state durability from replaceable repository-serving compute.

The distinction is central:
\begin{quote}
\textbf{Continuity addresses repository-state scalability; RaS investigates agent-state scalability.}
\end{quote}
Cursor Continuity does not test fresh-agent reconstruction, forced state reset, RRI, or RaS economics.

\subsection{Novelty boundary}

RaS does not claim novelty for Git, agents reading repositories, external memory, context retrieval, model routing, tests, durable execution, or stateless web services. The proposed research contribution is the combination of:

\begin{enumerate}
    \item repository-native durable engineering state as the primary continuity mechanism;
    \item deliberate disposal of high-capability reasoning state between engineering stages;
    \item Forced-State-Reset Evaluation;
    \item the proposed Repository Resumability Index;
    \item explicit reconstruction-cost measurement rather than treating rediscovery as free;
    \item separation of high-capability reasoning from stateful execution and privilege;
    \item investigation of agent-state scalability as distinct from repository-state scalability.
\end{enumerate}

Whether that combination proves practically important remains an empirical question.

\section{Threats to Validity}
\label{sec:threats}

RaS is vulnerable to several classes of confound that can make a forced-reset architecture appear stronger or weaker than it is.

\subsection{External validity}

\textbf{Single-repository bias} is immediate if early experiments use one codebase. Repository organisation, test quality, language, module boundaries, and engineering conventions can all affect reconstruction. \textbf{Single-model bias} is equally important because models differ in code navigation, long-context use, tool calling, and summarisation. \textbf{Repository complexity}, \textbf{task locality}, and \textbf{dependency width} may determine whether bounded reconstruction is feasible. Results in software engineering must not be silently generalised to personal assistants, embodied agents, research agents, or other domains.

\subsection{Internal validity}

\textbf{Task-selection bias} can favour RaS if experimenters choose tasks that are naturally encoded in source and tests. Conversely, selecting tasks that rely on deliberately undocumented intent could unfairly penalise it. \textbf{Historical commit survivor bias} arises when tasks are reconstructed only from changes known to have succeeded. \textbf{Future-history leakage} is especially serious: later tests, comments, files, or commit messages may reveal the answer to a historical task.

\textbf{Human intervention} can become an undeclared external memory channel. \textbf{Experiment-order effects} can occur if operators, task generators, or infrastructure learn across repeated conditions. \textbf{Contamination} can arise from model pretraining on public repositories or accidental sharing of artefacts between conditions. \textbf{Model nondeterminism} means a small number of runs can confuse stochastic variation with an architectural effect. \textbf{Provider updates} can alter behaviour during a study even when a model name is unchanged.

\subsection{Construct validity}

The project depends on several proxies. RRI measures correct continuation after reset, not the latent information equivalence expressed by the repository-sufficiency formulation. \textbf{Documentation quality} may dominate the result, causing the study to measure the benefit of unusually disciplined documentation rather than repositories generally. This is why semantic ablation is necessary.

\textbf{Reconstruction-cost classification} is another major risk. Repository reading that is necessary solely because history was destroyed must not disappear into a generic reasoning bucket. \textbf{Token telemetry gaps} can prevent complete cost comparison, particularly when providers expose caching or context accounting differently. \textbf{Prompt caching} may advantage persistent conditions or repeated repository prefixes in ways that must be measured rather than assumed.

\textbf{Hidden-verifier bias} can make success criteria too narrow, too aligned with the task author, or inadvertently informative. A hidden verifier should adjudicate behaviour, not encode an implementation preference.

\subsection{Reproducibility and statistical conclusion validity}

A private P0 subject limits independent reproduction. The appropriate response is not to publish proprietary material, but to classify the limitation explicitly and follow with public-repository replications. P0 will also have limited \textbf{statistical power}; it is intended to validate protocol and telemetry, not establish broad effect sizes.

Later studies must predeclare sample sizes, exclusions, missing-data handling, confidence intervals, and treatment of multiple comparisons. Failed or aborted runs should remain part of the experiment record unless an exclusion rule was committed before outcomes were observed.

\subsection{The central falsifier}

The most direct criticism is that RaS simply replaces conversation replay with repository rediscovery. If reconstruction tokens, file reads, tool calls, or wall-clock time grow sufficiently with repository size and depth, the architecture may be resumable but not economically useful. This possibility is not a secondary limitation. It is a central falsification condition for the state-economics claim.

\section{Discussion}
\label{sec:discussion}

RaS changes the locus of persistence. Instead of asking how to keep a reasoning process alive indefinitely, it asks what durable artefacts allow a fresh process to act correctly. That difference has implications for agent design, infrastructure, human oversight, and the interpretation of “memory” in software engineering.

\subsection{Repository state as semantic compression}

Engineering work routinely compresses large reasoning trajectories into smaller accepted artefacts. A debugging session can contain dozens of hypotheses and experiments, yet the durable consequence may be a five-line fix, one regression test, and a short design note. From a future-action perspective, this can be a desirable lossy compression: failed thoughts need not be replayed if the accepted state preserves the constraint that matters.

The important distinction is therefore between \emph{lossless conversational memory} and \emph{decision-relevant semantic state}. RaS predicts that the latter may be enough for many software-engineering continuations. Semantic-state ablation tests which artefacts perform that compression well and where important rationale is lost.

\subsection{Human and AI capability profiles}

The motivating workflow was asymmetric rather than purely autonomous. Human and AI systems possess different capability profiles. Human contribution may include persistent goals, grounded experience, problem selection, taste, risk judgement, strategic pivots, and deciding what evidence matters. High-capability AI may contribute broad technical synthesis, formalisation, rapid generation, critique, mathematical modelling, and implementation support. Execution agents contribute comparatively cheap or local machine interaction. Deterministic systems contribute compilation, tests, and measurements.

A useful conceptual workflow is:
\[
\text{human proposes/orchestrates}
\rightarrow
\text{high-capability AI amplifies/reasons}
\rightarrow
\text{execution worker implements}
\rightarrow
\text{tests/data/reality adjudicate}.
\]
The final arbiter is external evidence. This is discussion and observed motivation, not the central empirical RaS result.

\subsection{Fixed infrastructure and useful throughput}

If persistent reasoning-state burden per engineering programme can be reduced, fixed inference infrastructure may potentially support greater useful concurrency. Conceptually:
\[
\text{fixed infrastructure}
+
\text{lower persistent state requirement per programme}
\rightarrow
\text{potentially greater useful agent throughput}.
\]

This is an implication, not a measured result. Real serving throughput depends on model weights, KV-cache management, prompt sharing, batching, routing, hardware, request length, latency constraints, and provider architecture. Work such as vLLM demonstrates that KV-cache management materially affects serving throughput \citep{kwon2023vllm}, but it does not establish that RaS produces infrastructure savings.

The paper therefore makes no claim that datacentre constraints are solved, that a particular provider will save a given percentage, or that agent abundance follows on any specified timeline.

\subsection{What failure would teach}

A negative RaS result would still be valuable. If fresh agents fail because key rationale lives outside repositories, that identifies a missing state representation problem. If success depends heavily on concise state documents, the result may suggest that explicit semantic manifests are more important than source history. If reconstruction cost grows with dependency width, future systems may need better repository indexing or task-specific state compilers. If persistent agents outperform reset agents only on certain task classes, hybrid lifecycle policies become a more credible target than universal statelessness.

For that reason, the project should resist turning the architecture into a product claim before the experiments. The useful scientific question is not whether RaS can be made to look successful, but which engineering conditions make disposable reasoning viable.

\section{Repository-State and Agent-State Scalability}
\label{sec:scalability}

RaS separates two scaling problems that are easy to conflate.

The first is \emph{repository-state scalability}: how durable source history, objects, refs, and repository-serving infrastructure remain available as repository count and size grow. The second is \emph{agent-state scalability}: how reasoning processes preserve or reconstruct the information needed to continue useful work across long programmes.

Modern repository infrastructure provides an architectural precedent for the first problem. Cursor Continuity describes a write-ahead log persisted in S3-compatible object storage and local repository-serving state that can be materialised or rebuilt from the durable object state \citep{marti2026continuity}. The important RaS observation is not that this implementation proves anything about agents. It is that repository durability and repository-serving compute can be designed as different layers.

Figure~\ref{fig:scalability-stack} extends the separation conceptually.

\begin{figure}[htbp]
\centering
\input{figures/scalability-stack}
\caption{Conceptual separation of repository-state scalability and agent-state scalability. Cursor Continuity is relevant only to the upper repository-infrastructure precedent; the lower RaS layers remain hypotheses to be evaluated.}
\label{fig:scalability-stack}
\end{figure}

The conceptual stack is:

\begin{enumerate}
    \item durable object/log state;
    \item replaceable repository-serving compute;
    \item durable engineering repository state;
    \item replaceable high-capability reasoning;
    \item replaceable constrained execution workers.
\end{enumerate}

The top layers concern storage and Git serving. The bottom layers concern whether engineering semantics can be reconstructed after reasoning processes disappear. A repository system can scale perfectly while an agent still requires complete historical conversation. Conversely, an agent could be resumable from repository state even if the underlying Git service is implemented conventionally. The research questions are therefore independent.

This separation also clarifies the infrastructure implication. Long-context inference creates state-management pressure, including KV-cache allocation and bandwidth requirements; serving work such as PagedAttention shows that memory management can affect practical throughput \citep{kwon2023vllm}. RaS does not claim to remove those costs. Every reasoning transaction still requires inference state while it runs.

The defensible implication is conditional:
\[
\begin{array}{c}
\text{fixed inference infrastructure}\
+\
\text{lower persistent high-capability state requirement per useful programme}
\end{array}
\quad\Longrightarrow\quad
\text{potentially greater useful concurrency or throughput}.
\]

This remains an \textbf{IMPLICATION}, not an empirical result. Real throughput depends on model size, prompt length, batching, cache reuse, scheduler design, latency objectives, hardware, and provider implementation. RaS does not claim to solve data-centre scarcity, eliminate infrastructure constraints, guarantee abundance, or imply a particular labour-market outcome.

The relevant empirical quantity is more modest: how much high-capability input state and session persistence are required per successful engineering transition when repository reconstruction is included honestly. If those quantities decline without loss of engineering quality, the infrastructure implication becomes more plausible. If they do not, the implication should be rejected.

\section{Broader Authoritative-State Externalisation Principle}
\label{sec:authoritative-principle}

RaS may be a software-engineering-specific instance of a broader systems idea:

\begin{quote}
\textbf{Persist authoritative state. Reconstruct computation.}
\end{quote}

Suppose a domain has an authoritative state $A_t$, historical interaction $H_t$, a new task $U_t$, and a useful future action $X_{t+1}$. A broader externalisation hypothesis would ask whether:
\begin{equation}
P(X_{t+1}\mid H_t,A_t,U_t)
\approx
P(X_{t+1}\mid A_t,U_t).
\label{eq:authoritative-general}
\end{equation}

If Equation~\ref{eq:authoritative-general} held for a domain, long-term model-native history could have limited marginal value relative to reconstructing from authoritative domain state. Candidate stores might include CRM systems, transactional databases, ticket systems, event logs, document systems, or infrastructure state.

This is future work, not a result of the present paper. Software repositories have several properties that make them unusually favourable: versioning, explicit diffs, executable tests, deterministic build systems, inspectable schemas, mature access control, and a long engineering culture of encoding decisions into artefacts. Other domains may lack equivalent verification or may depend heavily on tacit, interpersonal, or continuously changing state.

There is also no reason to assume that a single authoritative store is sufficient. A production software programme may require repository state plus an issue tracker, deployment state, incident history, and external requirements. The broader principle should therefore be understood as an invitation to identify the \emph{authoritative domain state} rather than to force every agent into Git.

The proposed generalisation becomes scientifically useful only if it creates falsifiable questions: what state must survive, what computation can be safely reconstructed, how much reconstruction costs, and what quality is lost when model-native history is removed? RaS provides one concrete domain in which those questions can be asked experimentally.

\section{Future Work}
\label{sec:future}

The hostile audit concludes that P0 should proceed only after required protocol controls are frozen. The immediate next step is therefore corpus selection and preregistration under the audit gates, not execution.

P0 must first create a reproducible candidate task window, sequential-dependence rule, task hashes, canonical workspace materialisation procedure, FUTURE_HISTORY_LEAK_GATE, task-leakage audit, frozen hidden-verifier hashes, reconstruction-probe rubric, model/tool/environment manifest, stopping/rerun rules, and resource-classification rules.

After P0, repeated trials on the same complex repository are needed before any uncertainty estimate is meaningful. Public replications should then vary repository size, language, documentation quality, test quality, task locality, dependency width, and architecture.

Semantic-state ablation should remove small preregistered information classes rather than broadly degrading the repository. The scientific question is which durable artefacts have marginal continuity value.

Tiered execution belongs in a separate experiment because lower-cost exploration/execution is established prior work and changing executor capability would confound the core history treatment.

Cross-model continuity is particularly interesting but not part of P0. A design in which model family $M_A$ produces an accepted $R_t$ and a different family $M_B$ resumes from it without predecessor model-native state would provide stronger evidence that continuity resides in external project state rather than a model-specific session representation. It should be tested only after same-model causal identification works.

Reconstruction scaling should be tested directly. Synthetic irrelevant repository growth can hold task semantics constant; real tasks can vary dependency width. Plain search, symbol indexes, static-analysis graphs, repository maps, and model-guided retrieval can be compared without claiming any one is intrinsically RaS.

Security work should evaluate malicious repository instructions, tests, dependencies, commit metadata, poisoned durable documentation, work-package manipulation, executor compromise, and secret exposure. Least-privilege benefits remain architectural possibilities until adversarially measured.

The broader Authoritative-State Externalisation Principle remains speculative. Software repositories have unusually strong versioning, diffability, executable verification, and reproducible tooling; CRM systems, event logs, databases, and document stores may not offer equivalent semantics.

\section{Conclusion}
\label{sec:conclusion}

Repository-as-State asks a narrow systems question with potentially broad consequences for software-engineering agents: must durable engineering continuity live inside a persistent high-capability reasoning process?

The proposed alternative is to make the software project the persistent object. Durable repository state $R_t$ and current task $U_t$ are transformed by a bounded reconstruction function into the active reasoning view. High-capability reasoning proposes a transition. A stateful or constrained executor may implement it. Compilers, tests, runtime observations, and behavioural verification adjudicate the candidate. Only an accepted state becomes $R_{t+1}$. The high-capability reasoning process may then disappear.

RaS therefore does not primarily attempt to make high-quality reasoning cheap. It attempts to make high-quality reasoning disposable while preserving the engineering programme.

Paper v0.1 formalises that proposition through approximate repository sufficiency, a cumulative-context model under explicit assumptions, a provider-neutral cost decomposition, reconstruction cost and RTF, transactional reasoning, tiered execution, durable semantic state, selective documentation, least privilege, Forced-State-Reset Evaluation, the proposed Repository Resumability Index, and a separate state-reconstruction probe. It distinguishes repository-state scalability from agent-state scalability and places broader authoritative-state externalisation explicitly in future work.

The theory is intentionally falsifiable. Forced reset may reduce engineering quality. Reconstruction may become expensive as repositories or dependency width grow. Important rationale may remain outside repository state. Tiered execution may add more coordination cost than it saves. Security separation may introduce new confused-deputy risks. Any of these outcomes should narrow or reject the corresponding claim rather than be explained away after measurement.

The central systems principle remains:

\begin{quote}
\textbf{Persist authoritative state. Reconstruct computation.}
\end{quote}

Controlled empirical evaluation is not yet available. P0 has not been executed. The next step is an independent audit of the foundation and paper, followed by corpus selection and preregistration if the design survives that review.


\bibliographystyle{plainnat}
\bibliography{references}

\end{document}
